\chapter{Action, Symmetry, and Energy}

We are supposed to end with energy conservation, but the lecture does not begin there. Leonard Susskind first stops to repair a mistake from the previous meeting, then rebuilds the polar-coordinate example slowly enough that the method, the equations, and the physical interpretation all come into focus together. These notes follow that rhythm on purpose. They are based on Leonard Susskind's lecture, with selected board-frame evidence preserved through LazyingArt LLC and Video2Book as witnesses for equations, diagrams, and blackboard layout.

\section{Correcting the Polar-Coordinate Example}

The opening move is a correction. The term that appeared in the radial equation last time was called Coriolis; it should have been called centrifugal. But the real point is not the name. Susskind goes back to the problem of a particle moving in a plane and replays what he calls the rules of the game: choose coordinates, write down the kinetic and potential energies, form the Lagrangian, and let the Euler--Lagrange equations do the rest.

For a particle in two dimensions we could use \(x\) and \(y\), but for a central-force problem \(r\) and \(\theta\) are the natural coordinates. The radial component of velocity is simply
\begin{equation}
v_r=\dot r,
\end{equation}
while the transverse component is
\begin{equation}
v_\theta=r\dot\theta.
\end{equation}
This is the first place where the lecture insists on not being sloppy: \(\dot\theta\) is an angular velocity, but the actual linear speed perpendicular to the radius is \(r\dot\theta\).

Since these components are perpendicular, the kinetic energy is
\begin{equation}
T=\frac{m}{2}\left(v_r^2+v_\theta^2\right)
=\frac{m}{2}\left(\dot r^{\,2}+r^2\dot\theta^{\,2}\right).
\end{equation}
Now specialize to a potential that depends only on the distance from the origin:
\begin{equation}
U=U(r).
\end{equation}
Then the Lagrangian is
\begin{equation}
\mathcal{L}=\frac{m}{2}\left(\dot r^{\,2}+r^2\dot\theta^{\,2}\right)-U(r).
\end{equation}

Already one sees the next theme being planted. The fact that \(U\) depends on \(r\) but not on \(\theta\) is not merely a simplifying assumption; it is the first visible trace of rotational symmetry, and Susskind points to it early because he intends to come back to it later.

\begin{figure}[h]
\centering
\includegraphics[width=0.72\textwidth]{lecture_04_figure_02.png}

\begin{tikzpicture}[scale=1.2]
\draw[->] (-0.2,0) -- (2.2,0);
\draw[->] (0,-0.2) -- (0,2.0);
\draw[thick] (0,0) -- (1.55,1.05);
\draw (0.75,0) arc (0:34:0.75);
\node at (0.95,0.25) {$\theta$};
\end{tikzpicture}
\end{figure}

The board at this stage is valuable precisely because it shows the concrete example and the general formalism overlapping on the same blackboard. The polar-coordinate Lagrangian is still visible at the top, while underneath it Susskind has begun writing the generic derivative with respect to a generalized velocity. That is exactly how the lecture moves: not by abandoning the example, but by letting it open into general notation.

\section{Canonical Momenta and the Radial Equation}

The next step is the standard one. For each generalized coordinate \(q_i\), define the canonical momentum
\begin{equation}
\pi_i=\frac{\partial \mathcal{L}}{\partial \dot q_i}.
\end{equation}
The corresponding Euler--Lagrange equation is
\begin{equation}
\frac{d\pi_i}{dt}=\frac{\partial \mathcal{L}}{\partial q_i}.
\end{equation}
The blackboard is partly obscured at the crucial instant, so we should say this cautiously: the clean standard form above is not guessed from elsewhere, but it is the natural completion supported by both the transcript and the visible board sequence.

\begin{figure}[h]
\centering
\includegraphics[width=0.62\textwidth]{lecture_04_figure_03.png}
\end{figure}

Now return to the concrete variables \(r\) and \(\theta\). For the radial coordinate,
\begin{equation}
\pi_r=\frac{\partial \mathcal{L}}{\partial \dot r}=m\dot r.
\end{equation}
So
\begin{equation}
\frac{d\pi_r}{dt}=m\ddot r.
\end{equation}
On the right-hand side, \(r\) appears in two places: once inside \(U(r)\), and once multiplying \(\dot\theta^{\,2}\) in the kinetic term. Therefore
\begin{align}
\frac{\partial \mathcal{L}}{\partial r}
&=
\frac{\partial}{\partial r}
\left[
\frac{m}{2}r^2\dot\theta^{\,2}-U(r)
\right] \\
&=
mr\dot\theta^{\,2}-\frac{dU}{dr}.
\end{align}
The radial equation of motion is thus
\begin{equation}
m\ddot r=mr\dot\theta^{\,2}-\frac{dU}{dr}.
\end{equation}

This is the first payoff. The outward term
\begin{equation}
mr\dot\theta^{\,2}
\end{equation}
is not inserted by hand. It comes straight from differentiating the Lagrangian. This is one of Susskind's main practical claims: the equations of motion may look complicated, but the Lagrangian is often easy to write down. Once it is written down, the rest is in some sense algorithmic.

\subsection{Question \& Answer}

\paragraph{Question.} What does it mean to say that the centrifugal term is ``not a real force''?

\paragraph{Answer.} In the lecture Susskind immediately retracts the word ``real,'' and that is the right instinct. The useful distinction is narrower. The term \(mr\dot\theta^{\,2}\) does not come from differentiating a potential. It is not the \(-dU/dr\) piece. In that sense it is not a potential-derived force. But it is not fictitious in the weak sense of being ignorable. It appears on the right-hand side of the equation of motion, so it is part of the generalized radial force.

There is also a more physical way to say the same thing. If the particle has angular motion, then its velocity is not pointed straight at the origin. Left to itself it tends to continue in a straight line, and in polar variables that tendency appears as an outward contribution. The language can wobble; the equation does not.

\section{Rotational Symmetry and Angular Momentum}

Susskind now turns to the \(\theta\)-equation. This is where the example stops being merely a mechanics exercise and becomes a first instance of symmetry producing a conservation law.

The momentum conjugate to \(\theta\) is
\begin{equation}
\pi_\theta=\frac{\partial \mathcal{L}}{\partial \dot\theta}=mr^2\dot\theta.
\end{equation}
The Euler--Lagrange equation is
\begin{equation}
\frac{d\pi_\theta}{dt}=\frac{\partial \mathcal{L}}{\partial \theta}.
\end{equation}
But \(\theta\) does not appear in the Lagrangian at all. That is exactly the statement that the problem is rotationally invariant. Hence
\begin{equation}
\frac{\partial \mathcal{L}}{\partial \theta}=0,
\end{equation}
and so
\begin{equation}
\frac{d}{dt}\left(mr^2\dot\theta\right)=0.
\end{equation}
Therefore
\begin{equation}
mr^2\dot\theta=L,
\end{equation}
where \(L\) is constant. Susskind first calls this the canonical momentum conjugate to \(\theta\), and then reminds us that we already know its ordinary name:
\begin{equation}
L=\text{angular momentum}.
\end{equation}

This conservation law can now be fed back into the radial equation. Solve for \(\dot\theta\):
\begin{equation}
\dot\theta=\frac{L}{mr^2}.
\end{equation}
Substitute into the radial equation:
\begin{align}
m\ddot r
&=
-\frac{dU}{dr}
+
mr
\left(
\frac{L}{mr^2}
\right)^2 \\
&=
-\frac{dU}{dr}
+
\frac{L^2}{mr^3}.
\end{align}
So the radial dynamics can be rewritten as
\begin{equation}
m\ddot r=-\frac{dU}{dr}+\frac{L^2}{mr^3}.
\end{equation}

The lecture dwells on the meaning of this. Once the angular momentum is fixed by the initial conditions, the entire effect of the angular variable is condensed into an effective repulsive term proportional to \(1/r^3\). At small distances that is stronger than a \(1/r^2\) gravitational or Coulomb attraction. So a particle with nonzero angular momentum does not naturally plunge into the origin. This is the local physical lesson that Susskind wants us to feel before he turns to the general theorem.

\section{Variation Notation as Preparation}

At this point the lecture pauses. Susskind says it is ``not really mathematics,'' only notation that we will see over and over again. But the pause is deliberate: it prepares the ground for the variational proof that follows.

Suppose \(f\) depends on many variables \(\alpha_i\). For small changes \(\delta \alpha_i\), the first-order change in \(f\) is
\begin{equation}
\delta f=\sum_i \frac{\partial f}{\partial \alpha_i}\,\delta\alpha_i.
\end{equation}
If \(\delta f=0\) for arbitrary variations \(\delta\alpha_i\), then necessarily
\begin{equation}
\frac{\partial f}{\partial \alpha_i}=0
\qquad
\text{for all } i.
\end{equation}
In the lecture this is described in the language of a minimum. On the page it is cleaner to say that the function is stationary to first order. The notation \(\delta f=0\) is simply shorthand for the vanishing of all first derivatives.

Exactly the same notation is then applied to the action:
\begin{equation}
A=\int \mathcal{L}\,dt,
\qquad
\delta A=0.
\end{equation}
Susskind writes the action as \(A\), not \(S\), and since that is also what the blackboard shows, it is best to keep the same notation here.

\begin{figure}[h]
\centering
\includegraphics[width=0.72\textwidth]{lecture_04_figure_04.png}

\begin{tikzpicture}[scale=0.95]
\draw[->] (0,0) -- (0,4.2);
\draw[->] (0,0) -- (3.8,0);
\draw[thick] (1.0,0.35) .. controls (1.35,1.2) and (0.95,2.45) .. (1.65,3.8);
\draw[thick,dashed] (1.32,0.5) .. controls (1.67,1.35) and (1.25,2.6) .. (1.98,3.95);
\draw (1.10,0.95) -- (1.43,1.05);
\draw (1.18,1.75) -- (1.51,1.85);
\draw (1.20,2.55) -- (1.54,2.65);
\draw (1.42,3.25) -- (1.76,3.35);
\node[below] at (2.0,0) {$q$};
\node[left] at (0,3.9) {$t$};
\end{tikzpicture}
\end{figure}

In clean notation, the first variation of the action is
\begin{equation}
\delta A
=
\int
\sum_i
\left(
\frac{\partial \mathcal{L}}{\partial q_i}\,\delta q_i
+
\frac{\partial \mathcal{L}}{\partial \dot q_i}\,\delta \dot q_i
\right)dt.
\end{equation}
The board does not show every symbol sharply enough to justify a literal transcription, so this formula should be read as the cautious standard reconstruction that matches both the transcript and the visible layout of the blackboard.

\section{Infinitesimal Symmetries and Noether's Argument}

Now the lecture turns to what Susskind explicitly calls the heart of the evening: the connection between symmetries and conservation laws.

A symmetry is a change one can make in the system without changing the action. For present purposes we are interested in infinitesimal symmetries, that is, transformations built from arbitrarily small displacements:
\begin{equation}
\delta q_i=\epsilon f_i(q).
\end{equation}
The functions \(f_i\) may depend on all the coordinates. This already includes translations and rotations, and it leaves room for more general transformations as well.

At this point the lecture briefly detours through a garbled discussion of a die or cube. The clean conceptual point is simple: discrete symmetries are not the same thing as the infinitesimal symmetries used here. A sphere is invariant under arbitrarily small rotations; a cube is not. Noether's argument, in the form developed here, is about the continuous case.

Now take a trajectory that solves the equations of motion between two times \(t_1\) and \(t_2\). For an ordinary least-action variation, the endpoints are held fixed. But a symmetry-generated variation is different: it moves the entire trajectory, including its endpoints. That is exactly why the proof feels slightly different from the usual least-action derivation. The vanishing of \(\delta A\) will now come from symmetry, not from the restriction to endpoint-fixed variations.

After integrating by parts, the variation of the action takes the standard form
\begin{equation}
\delta A
=
\int
\sum_i
\left[
\frac{\partial \mathcal{L}}{\partial q_i}
-
\frac{d}{dt}
\left(
\frac{\partial \mathcal{L}}{\partial \dot q_i}
\right)
\right]
\delta q_i\,dt
+
\left[
\sum_i
\frac{\partial \mathcal{L}}{\partial \dot q_i}\,\delta q_i
\right]_{t_1}^{t_2}.
\end{equation}
Again, this is the cautious standard reconstruction of the board formula; the blackboard shows the structure clearly, even if not every symbol is fully legible.

\begin{figure}[h]
\centering
\includegraphics[width=0.78\textwidth]{lecture_04_figure_05.png}

\begin{tikzpicture}[scale=1.0]
\draw[->] (0,0) -- (0,4.4);
\draw[->] (0,0) -- (4.1,0);
\draw[thick] (1.05,0.35) .. controls (0.90,1.35) and (1.85,2.35) .. (1.45,3.95);
\draw[thick,dashed] (1.42,0.52) .. controls (1.27,1.52) and (2.18,2.52) .. (1.82,4.10);
\draw (1.02,0.90) -- (1.38,1.02);
\draw (1.01,1.70) -- (1.37,1.82);
\draw (1.15,2.50) -- (1.52,2.62);
\draw (1.32,3.28) -- (1.67,3.40);
\node[below] at (2.0,0) {$q$};
\node[left] at (0,4.0) {$t$};
\end{tikzpicture}
\end{figure}

Now the lecture's logic becomes very economical. Because the original path is a solution of the equations of motion, the bulk term vanishes:
\begin{equation}
\frac{\partial \mathcal{L}}{\partial q_i}
-
\frac{d}{dt}
\left(
\frac{\partial \mathcal{L}}{\partial \dot q_i}
\right)
=0.
\end{equation}
Therefore
\begin{equation}
\delta A=
\left[
\sum_i
\frac{\partial \mathcal{L}}{\partial \dot q_i}\,\delta q_i
\right]_{t_1}^{t_2}.
\end{equation}
But the variation has been chosen to be a symmetry, so \(\delta A=0\) for that reason. Hence
\begin{equation}
\left[
\sum_i
\frac{\partial \mathcal{L}}{\partial \dot q_i}\,\delta q_i
\right]_{t_1}^{t_2}=0.
\end{equation}
Now insert \(\delta q_i=\epsilon f_i(q)\):
\begin{equation}
\epsilon
\left[
\sum_i
\frac{\partial \mathcal{L}}{\partial \dot q_i}\,f_i(q)
\right]_{t_1}^{t_2}=0.
\end{equation}
Dropping the overall factor \(\epsilon\), and writing
\begin{equation}
\pi_i=\frac{\partial \mathcal{L}}{\partial \dot q_i},
\end{equation}
we obtain
\begin{equation}
\left[
\sum_i \pi_i f_i(q)
\right]_{t_1}^{t_2}=0.
\end{equation}
So the quantity
\begin{equation}
Q=\sum_i \pi_i f_i(q)
\end{equation}
has the same value at any two times. This is the conserved Noether charge associated with the symmetry.

The lecture pauses here to note two things. First, the sum over \(i\) is always present, even when it is not written out every time in speech. Second, the word ``charge'' is just physics slang for ``conserved quantity.'' Momentum is a charge in this sense, angular momentum is a charge, and energy will be one too.

\section{Worked Noether Charges: Translation and Rotation}

Once the general argument is on the board, Susskind immediately says: let us work out an example or two. This is important pedagogically. The theorem is not left floating in abstraction; it is checked against the familiar conservation laws we already know.

For translation along the \(x\)-direction, every particle is shifted by the same amount:
\begin{equation}
\delta x_i=\epsilon,
\qquad
\delta y_i=\delta z_i=0.
\end{equation}
So the generator is \(f_{x_i}=1\) and all the other components vanish. The Noether charge is
\begin{equation}
Q_x=\sum_i \pi_{x_i}.
\end{equation}
For the usual mechanical Lagrangian, \(\pi_{x_i}=m_i\dot x_i=p_{x_i}\), and therefore
\begin{equation}
Q_x=\sum_i p_{x_i},
\end{equation}
the total \(x\)-component of momentum.

Again the lecture stresses the remarkable economy of the result. We did not need to know the detailed form of the internal forces. We needed only that the action of the isolated system is unchanged when the whole system is translated.

Now consider an infinitesimal rotation in the \(x\)-\(y\) plane. The lecture wobbles briefly on the sign, then settles on the correct first-order formulas:
\begin{equation}
\delta x=-\epsilon y,
\qquad
\delta y=\epsilon x.
\end{equation}
Hence
\begin{equation}
f_x=-y,
\qquad
f_y=x.
\end{equation}
The corresponding Noether charge is
\begin{align}
Q_{\mathrm{rot}}
&=\pi_x f_x+\pi_y f_y \\
&=-\pi_x y+\pi_y x \\
&=xp_y-yp_x.
\end{align}
This is the angular momentum perpendicular to the plane. For many particles, one sums the same expression over the particles.

So the theorem reproduces the familiar conservation laws in exactly the way Susskind promised. Translation invariance gives momentum. Rotational invariance gives angular momentum. The theorem does not manufacture strange new objects first; it explains why the old ones were conserved all along.

\section{Time Translation, the Hamiltonian, and Energy}

At this stage the lecture turns to the promised destination. If the laws of the system have no explicit time dependence, then a solution shifted forward in time is still a solution. This is time-translation invariance.

Here the geometry is a little subtler than before. One can think of the whole trajectory as being pushed upward in the \(t\)-direction. But at a fixed time \(t\), this also appears as a horizontal displacement in configuration space, and that horizontal shift depends on the slope of the trajectory at the point in question.

\begin{figure}[h]
\centering
\begin{tikzpicture}[scale=1.0]
\draw[->] (0,0) -- (0,5.0);
\draw[->] (0,0) -- (4.6,0);
\draw[thick] (1.55,0.55) .. controls (1.15,1.45) and (2.20,2.70) .. (1.70,4.15);
\draw[thick,dashed] (1.20,1.10) .. controls (0.82,2.00) and (1.88,3.25) .. (1.36,4.70);
\draw[dotted] (1.55,0.55) -- (1.20,1.10);
\draw[dotted] (1.70,4.15) -- (1.36,4.70);
\node[below] at (2.3,0) {$q$};
\node[left] at (0,4.7) {$t$};
\node[right] at (1.78,4.15) {$t_2$};
\node[right] at (1.62,0.55) {$t_1$};
\end{tikzpicture}
\end{figure}

The transformed trajectory may be written as
\begin{equation}
q_i(t)\mapsto q_i(t-\epsilon).
\end{equation}
Expanding to first order gives
\begin{equation}
q_i(t-\epsilon)=q_i(t)-\epsilon \dot q_i(t)+O(\epsilon^2),
\end{equation}
so the induced variation is
\begin{equation}
\delta q_i=-\epsilon \dot q_i.
\end{equation}

\subsection{Question \& Answer}

\paragraph{Question.} Why is there a minus sign in \(\delta q_i=-\epsilon \dot q_i\)?

\paragraph{Answer.} Because after we push the whole trajectory upward in time, the point on the new curve at fixed time \(t\) comes from the old curve at the earlier time \(t-\epsilon\). If the trajectory slopes to the right, then moving it upward shifts the corresponding point to the left. Susskind explicitly says that he always gets confused by this sign and has to draw the picture again; the page should keep that honesty. The sign is not decorative. It is a real geometric point.

\paragraph{Question.} Why is there now an extra \(\mathcal{L}\)-term in the Noether argument?

\paragraph{Answer.} Because a time shift changes the time interval on which the action is evaluated. The shifted trajectory does not begin and end at the same times as the original one. There are therefore two small overhang regions: one near the lower end, one near the upper end. Each contributes a small piece of action. For sufficiently small \(\epsilon\),
\begin{equation}
A_{\text{overhang}}\approx \mathcal{L}(t_2)\,\epsilon,
\qquad
B_{\text{overhang}}\approx \mathcal{L}(t_1)\,\epsilon.
\end{equation}
This is the new ingredient. The ordinary Noether endpoint term is still present, but now it is accompanied by these extra pieces because the integration interval itself has shifted.

We can now run the argument. Compare the original trajectory and the time-shifted one over their common interval \([t_1,t_2]\), and then add the overhang contributions. The bulk Euler--Lagrange term vanishes because the original path satisfies the equations of motion. What remains is
\begin{equation}
0
=
\left[
\sum_i
\frac{\partial \mathcal{L}}{\partial \dot q_i}\,\delta q_i
\right]_{t_1}^{t_2}
+
\mathcal{L}(t_2)\epsilon
-
\mathcal{L}(t_1)\epsilon.
\end{equation}
Now substitute \(\delta q_i=-\epsilon \dot q_i\):
\begin{equation}
0
=
\epsilon
\left[
\mathcal{L}
-
\sum_i
\dot q_i
\frac{\partial \mathcal{L}}{\partial \dot q_i}
\right]_{t_1}^{t_2}.
\end{equation}
Hence the quantity
\begin{equation}
\mathcal{L}
-
\sum_i
\dot q_i
\frac{\partial \mathcal{L}}{\partial \dot q_i}
\end{equation}
is conserved.

The lecture then makes a small but important sign choice. Susskind first identifies the conserved quantity in exactly the form above, then immediately says that if a quantity is conserved, so is its negative. He therefore flips the sign and defines
\begin{equation}
H=\sum_i \dot q_i\,\pi_i-\mathcal{L}.
\end{equation}
This is the Hamiltonian. In the language of the lecture, it is also the energy.

To check that the definition agrees with the familiar expression, take
\begin{equation}
\mathcal{L}=\frac{1}{2}m\dot x^2-U(x).
\end{equation}
Then
\begin{equation}
\pi_x=\frac{\partial \mathcal{L}}{\partial \dot x}=m\dot x,
\end{equation}
and so
\begin{align}
H
&=
\pi_x\dot x-\mathcal{L} \\
&=
m\dot x^2-\left(\frac{1}{2}m\dot x^2-U(x)\right) \\
&=
\frac{1}{2}m\dot x^2+U(x).
\end{align}
Thus the Hamiltonian becomes the familiar total energy \(T+U\).

This is the point at which the lecture deliberately becomes more fundamental. Energy is not merely the combination \(T+U\) that we happened to know in elementary mechanics. More deeply, energy is the conserved quantity associated with time-translation invariance. If the Lagrangian has the ordinary form \(T-U\), that conserved quantity happens to be \(T+U\). But the Noether derivation is the deeper statement.

There is also a sharp caveat. If the Lagrangian depends explicitly on time,
\begin{equation}
\mathcal{L}=\mathcal{L}(q,\dot q,t),
\end{equation}
then the time-shifted path need not have the same action. In that case the symmetry is lost, and so is the conservation law. Energy need not be conserved.

The lecture closes by stressing that this is the most general definition in the present framework: the Lagrangian is the integrand in the action, and when time-translation symmetry is present, the associated Noether charge is the energy. The next step is then natural. One has been led directly to the Hamiltonian, so Hamiltonian mechanics is what comes next.

\section{Summary}

The lecture begins with a repair to a familiar mechanics problem and ends with the general definition of energy. That progression is essential. First we see, in polar coordinates, how the Euler--Lagrange equations produce the centrifugal term and how rotational symmetry produces angular-momentum conservation. Then the notation of variations is put in order. Only after that does the lecture turn to the general symmetry argument and the Noether charge
\begin{equation}
Q=\sum_i \pi_i f_i(q).
\end{equation}

Time translation is the final and special case. Because the time interval itself shifts, the ordinary Noether endpoint term is supplemented by the overhang contributions, and together they yield the conserved Hamiltonian
\begin{equation}
H=\sum_i \dot q_i\,\pi_i-\mathcal{L}.
\end{equation}
For ordinary mechanical systems this reduces to \(T+U\), but the lecture's deeper claim is that the Hamiltonian is the more fundamental notion: it is the quantity conserved because the laws do not change when the experiment is begun a little earlier or a little later. From there the road leads directly to Hamiltonian mechanics.